\chapter{Appendix 03: Gamma Kitchen}

\section{The Legacy Paradigm \\\& Its Failures}
Legacy food systems depend on gig-economy delivery apps characterized by dark patterns, addictive notifications, and opaque algorithms. They require constant attention and cloud connectivity, ignoring the communal and physical aspects of food preparation while draining device batteries.

\section{The Incremental Automation Pathway}
Phase 1 shifts inventory management to local CRDT ledgers shared among neighbors. Phase 2 introduces ambient interfaces, utilizing low-power E-ink displays on pantries to show ingredient availability. Phase 3 automates communal meal planning, with the system autonomously sourcing local ingredients based on thermodynamic energy constraints.

\section{The Human Boundary (Limits of Automation)}
Food is inherently cultural and personal. The system will never automate the final recipe selection or the physical cooking process without human initiation. The interface steps back during meal preparation, deliberately minimizing cognitive load and allowing humans to focus on the tactile experience.

\section{Interface Mechanics (The "How")}
Spatial UI is used sparingly to project precise measurements onto cutting boards. The system degrades UI richness during high-energy cooking (e.g., using ovens) to save power. We identified a gap: an ambient non-visual cue system (like localized auditory nudges) for cooking times, reducing screen reliance entirely.
